\section{System Implementation}

\subsection{Semaphore and Synchronization Strategy}

The system uses a producer-consumer architecture with multiple semaphores and queues to synchronize concurrent RTOS
tasks.

The Sensor Monitoring Task acts as the producer. Every 5 seconds, it reads the latest temperature and humidity values
from the DHT20 sensor and stores them in a shared queue. After new data is available, the task releases the semaphore
\texttt{sensor\_ready\_sem} to notify the Decision Maker Task.

The Decision Maker Task acts as the central coordinator. It waits for \texttt{sensor\_ready\_sem}, retrieves the newest
sensor reading, and distributes the data to the LCD and Heater tasks through dedicated queues. Depending on the
measured temperature and humidity, it may also activate the Cooler Task or Humidifier Task by releasing the
corresponding semaphores.

\begin{center}
\renewcommand{\arraystretch}{1.25}

\begin{tabular}{|l|l|}
    \hline
    \rowcolor{RTOSBlue}
    \color{white}\bfseries\sffamily Semaphore / Queue &
    \color{white}\bfseries\sffamily Purpose \\
    \hline

    \texttt{sensor\_ready\_sem} &
    Signals that a new sensor sample is available. \\
    \hline

    \texttt{lcd\_sem} &
    Signals that LCD data are ready to display. \\
    \hline

    \texttt{heater\_sem} &
    Activates the heater task. \\
    \hline

    \texttt{cooler\_sem} &
    Activates the cooler task. \\
    \hline

    \texttt{humidifier\_sem} &
    Activates the humidifier task. \\
    \hline

    \texttt{cooler\_finished\_sem} &
    Signals that the cooler has completed its cycle. \\
    \hline

    \texttt{humidifier\_finished\_sem} &
    Signals that the humidifier has completed its sequence. \\
    \hline
\end{tabular}

\end{center}

This design prevents race conditions between tasks and ensures that control actions are executed only when valid sensor
data is available. Furthermore, the completion semaphores guarantee that the Decision Maker waits for actuator
operations to finish before continuing with further control decisions.

\subsection{Read Temperature Task}

\textbf{Priority:} Medium

\textbf{Period:} 5 seconds

The Read Temperature Task periodically acquires temperature and humidity measurements from the DHT20 sensor. The
collected data is stored in a shared queue and signaled through the \texttt{sensor\_ready\_sem} semaphore. This task
acts as the producer in the producer-consumer communication pattern.

\textbf{Main implementation:}

\begin{lstlisting}
async def task_monitoring():
    while True:
        data = {
            "temp": await dht20.atemperature(),
            "humi": await dht20.ahumidity()
        }
        add_latest(sensor_queue, sensor_ready_sem, data)
        await asleep_ms(5000)
\end{lstlisting}

The task continuously generates environmental measurements and provides them to the Decision Maker Task for further
processing.

\subsection{Cooler Task}

\textbf{Priority:} High

\textbf{Trigger Condition:} Temperature \textgreater{} 28°C

The Cooler Task is event-driven and remains blocked until the \texttt{cooler\_sem} semaphore is released by the
Decision Maker Task. Once activated, the cooler LED turns GREEN for a fixed duration of 5 seconds, representing the
cooling process. After the timer expires, the LED is turned OFF and a completion signal is sent through
\texttt{cooler\_finished\_sem}.

\textbf{Main implementation:}

\begin{lstlisting}
async def task_cooler():
    while True:
        await cooler_sem.acquire()
        rgb_led_D5.show(0, hex_to_rgb('#00ff00'))
        await asleep_ms(5000)
        rgb_led_D5.show(0, hex_to_rgb('#000000'))
        cooler_finished_sem.release()
\end{lstlisting}

\textbf{State Machine:}

\begin{itemize}
  \item IDLE (BLACK): Cooler inactive.
  \item COOLING (GREEN): Cooler active for 5 seconds.
  \item Return to IDLE after timer expiration.
\end{itemize}

The Cooler Task state machine is shown in Figure~\ref{fig:cooler-state-machine}.

% \begin{figure}[htbp]
%   \centering
%   \includegraphics[width=0.90\textwidth]{images/cooler_state_machine.png}
%   \caption{Cooler control state machine.}
%   \label{fig:cooler-state-machine}
% \end{figure}

\subsection{Heater Task}

\textbf{Priority:} Medium

\textbf{Trigger Condition:} New temperature reading available.

The Heater Task provides a visual indication of environmental temperature conditions. The task waits for a signal from
\texttt{heater\_sem}, retrieves the latest temperature value, and updates the RGB LED according to predefined
temperature ranges.

\textbf{Main implementation:}

\begin{lstlisting}
async def task_heater():
    while True:
        await heater_sem.acquire()
        temp = heater_queue.pop(0)["temp"]
        if 20 <= temp <= 25:
            rgb_led_D3.show(0, hex_to_rgb('#00ff00'))
        elif 15 <= temp <= 30:
            rgb_led_D3.show(0, hex_to_rgb('#ffff00'))
        else:
            rgb_led_D3.show(0, hex_to_rgb('#ff0000'))
\end{lstlisting}

\textbf{State Machine:}

\begin{itemize}
  \item SAFE (GREEN): 20°C ≤ Temp ≤ 25°C
  \item WARNING (YELLOW): 15°C ≤ Temp \textless{} 20°C or 25°C \textless{} Temp ≤ 30°C
  \item CRITICAL (RED): Temp \textless{} 15°C or Temp \textgreater{} 30°C
\end{itemize}

The Heater Task continuously reflects the current thermal condition of the environment.

\subsection{Humidifier Task}

\textbf{Priority:} High

\textbf{Trigger Condition:} Humidity \textless{} 40\%

The Humidifier Task is activated when the humidity falls below the predefined threshold. After receiving a signal
through \texttt{humidifier\_sem}, it executes a timed three-stage sequence representing the humidification process.

\textbf{Main implementation:}

\begin{lstlisting}
async def task_humidifier():
    while True:
        await humidifier_sem.acquire()
        rgb_led_D7.show(0, hex_to_rgb('#00ff00'))
        await asleep_ms(5000)
        rgb_led_D7.show(0, hex_to_rgb('#ffff00'))
        await asleep_ms(3000)
        rgb_led_D7.show(0, hex_to_rgb('#ff0000'))
        await asleep_ms(2000)
        rgb_led_D7.show(0, hex_to_rgb('#000000'))
        humidifier_finished_sem.release()
\end{lstlisting}

\textbf{State Machine:}

\begin{itemize}
  \item IDLE (BLACK): Humidifier inactive.
  \item GREEN: Active humidification for 5 seconds.
  \item YELLOW: Intermediate phase for 3 seconds.
  \item RED: Final phase for 2 seconds.
  \item Return to IDLE and re-check humidity.
\end{itemize}

The completion semaphore informs the Decision Maker that the humidification cycle has finished.
